\documentclass[10pt,a4paper]{amsart}
\usepackage[margin=19mm]{geometry}
\usepackage{amsmath,amssymb,mathtools}
\usepackage[scr=rsfs,cal=euler]{mathalfa}
\usepackage{xfrac}
\usepackage[unicode,hidelinks,bookmarks=false]{hyperref}
\newcommand{\ie}{i.e.\ }
\newcommand{\cf}{cf.\ }
\newcommand{\resp}{respectively\ }
\newcommand{\Subsection}{\subsection}
\providecommand{\nobreakdash}{\nobreak\hbox{-}}
\setlength{\emergencystretch}{2em}
\multlinegap0pt
\usepackage{amssymb}
\usepackage{caption}
\usepackage[shortlabels]{enumitem}
\usepackage{booktabs}
\usepackage{algorithm}\usepackage{algpseudocode}
\usepackage{tikz}
\usepackage{multirow}
\newtheorem{lemma}{Lemma}
\newcommand{\AYr}[1]{{#1}}
\newcommand{\QGr}[1]{{#1}}
\title{Oriented discrepancy of Hamilton cycles in oriented graphs satisfying Ore-type condition}
\author{Jiangdong Ai, Qiwen Guo, Gregory Gutin, Yongxin Lan, Qi Shao, Anders Yeo, Yacong Zhou}
\date{Source: arXiv:2501.05968v2 (CC BY 4.0)}
\begin{document}
\maketitle

\noindent\emph{Source and licence.} Oriented discrepancy of Hamilton cycles in oriented graphs satisfying Ore-type condition. Version arXiv:2501.05968v2 (CC BY 4.0), distributed by arXiv under the Creative Commons Attribution 4.0 International licence, http://creativecommons.org/licenses/by/4.0/, as recorded in the arXiv licence metadata for this exact version and shown on the abstract page https://arxiv.org/abs/2501.05968v2. Full licence notice: \texttt{licenses/arxiv-2501.05968-CC-BY-4.0.txt}.\par\smallskip

\noindent\textit{Editorial scope.} Counted unit: the article's lemma pair (source0.tex lines 397--401). The article's complete original proof of it is delivered with it (source0.tex lines 403--421, one \texttt{proof} environment of 19 lines). No mathematics is written, repaired or re-derived: the statement and the proof are the article's own lines, in the article's own notation, and no retained line is substituted or re-flowed.  No further source-local context is needed: every cross-reference of every retained line resolves inside the two delivered spans.  Nothing was omitted from any retained span, so the package carries no omission record.  No retained line contains a citation, so the wrapper carries no bibliography and no bibliography entry.  No retained line contains a figure environment, an image-inclusion command or drawing code, so no image is delivered and the package declares no asset dependency.  Editorial formatting only, in addition: an AMS \texttt{amsart} preamble that restates the article's own declarations and macros used by the retained lines, namely the environments \texttt{lemma} on separate counters, together with the standard packages those lines need.  The retained environments print in this document as Lemma 1 in order of appearance, whereas the article prints the same items with the numbering of its own full document.  The complete native source of the article as distributed in the arXiv e-print for this exact version is bundled unchanged as \texttt{sources/171445/source0.tex}.
	\begin{lemma}\label{lem:add x between pair}
	% Let $P = u_1u_2\ldots u_n$ be an oriented path and let $\{v_1,v_2,\dots,v_k\}$ be a set of vertices that are not in $V(P)$, such that $v_i$ is adjacent to $u_j$ for each $i\in [k]$ and $j\in [n]$. If $n\ge k+2$, then there exists an oriented path $P'$ with $\sigma^+(P') \ge \sigma^+(P)$, which is obtained from $P$ by adding all $v_i$ between $u_{a_i}$ and $u_{a_i+1}$ for $a_i\in [n-1]$, where $a_i\ne a_j$ for $i\ne j$. Similarly, an oriented path $P''$ satisfying $\sigma^-(P'') \ge \sigma^-(P)$ also exists.

    \QGr{Let $D$ be an oriented graph which contains an oriented path $P$ and a set of vertices $K$ disjoint from $V(P)$, such that $V(P)\subseteq N(v)$ for each $v\in K$. If $|V(P)|\ge |K|+2$, then $D$ contains an oriented path $P'$ with $V(P')=V(P)\cup K$ and same initial and terminal vertices as $P$, satisfying $\sigma^+(P')\ge \sigma^+(P) $. Similarly, an oriented path $P''$ satisfying $\sigma^-(P'') \ge \sigma^-(P)$ also exists.}
	\end{lemma}

    \begin{proof}
		% Note that it suffices to prove the existence of $P'$ when \QGr{$|V(P)|=|K|+2$.}  \QGr{Let $K=\{v_1,v_2,\dots,v_k\}$ and  $P = u_1u_2\ldots u_{k+2}$.} For $k=1$, if the arc between $v_1$ and $u_2$ is $v_1u_2$, then $P' = u_1v_1u_2u_3$ satisfies $\sigma^+(P') \ge \sigma^+(P)$. Otherwise, the arc between $v_1$ and $u_2$ is $u_2v_1$, and then $P' = u_1u_2v_1u_3$ satisfies $\sigma^+(P') \ge \sigma^+(P)$. We now prove that for $k\ge 2$, the oriented path $P'$ exists by induction on $k$.
		
		% For any $t\in [k+1]$, let \QGr{$P_{t} = u_1u_2\ldots u_{t}$.} Note that $\sigma^+(P_{t})\ge \sigma^+(P)-(k+2-t)$. Suppose to the contrary that such an oriented path $P'$ does not exist for $k$. Then, for every $i\in [k]$, the arc $v_iu_{k+1}$ exists. Otherwise, the arc between $v_{i_0}$ and $u_{k+1}$ is $u_{k+1}v_{i_0}$ for some $i_0\in [k]$. By the induction hypothesis, there is an oriented path $P'_{k+1}$ with \QGr{$V(P_{k+1}')=V(P_{k+1})\cup (K\setminus \{v_{i_0}\})$ and same initial and terminal vertices as $P_{k+1}$, satisfying $\sigma^+(P_{k+1}')\ge \sigma^+(P_{k+1}) $.} Let $P'$ be the oriented path which consists of $P'_{k+1}$, the arc $u_{k+1}v_{i_0}$ and the arc between $v_{i_0}$ and $u_{k+2}$. We have $\sigma^+(P')\ge \sigma^+(P'_{k+1}) +1\ge \sigma^+(P)$, which is a contradiction. Now we claim that \QGr{arc $v_iu_{t}$ exists for every $i\in [k]$ and some $t\in [k]$ if arc $v_iu_{j}$ exists for every $i\in [k]$ and $t+1 \le j\le k+1$. If this claim is not true, then the arc between $v_{i_0}$ and $u_t$ is $u_tv_{i_0}$ for some $i_0\in [k]$. We arbitrarily choose $v_{i_1}, \dots, v_{i_{k+1-t}}$ for $i_1,\dots,i_{k+1-t}\in [k]\setminus\{i_0\}$. By induction hypothesis, there is an oriented path $P'_{t}$ with $V(P_{t}')=V(P_{t})\cup (K\setminus \{v_{i_0},v_{i_1}, \dots, v_{i_{k+1-t}}\})$ and same initial and terminal vertices as $P_{t}$, satisfying $\sigma^+(P_{t}')\ge \sigma^+(P_{t})$. Let $P'$ be the oriented path which consists of $P'_{t}$, the directed path $u_tv_{i_0}u_{t+1}$, arcs $v_{i_j}u_{t+j}$ and $v_{i_j}u_{t+j+1}$ for each $j\in [k-t]$, the arc $v_{i_{k+1-t}}u_{k+1}$ and the arc between $v_{i_{k+1-t}}$ and $u_{k+2}$. We have $\sigma^+(P')\ge \sigma^+(P'_t) +2+ k-t\ge \sigma^+(P)$, which is a contradiction. By this claim, since the arc $v_iu_{k+1}$ exist for every $i\in [k]$, arcs $v_iu_j$ exists for every $i\in [k]$ and $j\in [k+1]$. Let $P'=u_1v_1u_2v_2u_3\dots u_kv_ku_{k+1}u_{k+2}$ be the oriented path which consists of arcs $v_iu_i$ and $v_iu_{i+1}$ for each $i\in [k]$ and the arc between $u_{k+1}$ and $u_{k+2}$.} Then $P'$ has the same initial and terminal vertices as $P$ and satisfies $V(P')=V(P)\cup K$ and $\sigma^+(P')\ge \sigma^+(P)$, which is a contradiction.
        \AYr{  Let $K=\{v_1,v_2,\dots,v_k\}$ and  $P = u_1u_2\ldots u_{l}$, where $l \geq k+2$. Let $A=\{a_1,a_2,\ldots,a_k\}$ be $k$ distinct arcs chosen randomly  and uniformly from $A(P)$. For each $i \in [k]$ let $a_i=s_it_i$ and let $P'$ be obtained from the path $P$,
        by substituting $a_i$ with the oriented path $s_i v_i t_i$ for each $i \in [k]$.
        
          Let $q$ and $r$ be any integers such that $2 \leq q \leq l-1$ and $r \in [k]$.
          If  $u_q v_r \in A(D)$ then $u_qv_r$ will be a forward arc in $P'$ with probability $\frac{1}{l-1}$, as this is the case if and only if $a_r = u_q u_{q+1}$.
        And, if  $v_r u_q \in A(D)$ then $v_r u_q$ will be a forward arc in $P'$ with probability $\frac{1}{l-1}$, as
         this is the case if $a_r = u_{q-1} u_{q}$. So every arc between $K$ and $\{u_2,u_3,u_4, \ldots , u_{l-1}\}$ will be a forward arc
        with probability $\frac{1}{l-1}$ and there are $k \cdot (l-2)$ such arcs.  And we remove at most $k$ forward arcs from $P$ when we remove the arcs in $A$. So, on average, $P'$
        will satisfy the following (as $l \geq k+2$),
         \[
        \sigma^+(P') \geq \sigma^+(P) + \frac{k\cdot (l-2)}{l-1} - k \geq \sigma^+(P) + \frac{k\cdot k}{k+1} - k >  \sigma^+(P) -1
        \]
        As $\sigma^+(P')$ is an integer this implies that there exists an oriented path $P'$ with $\sigma^+(P') \geq \sigma^+(P)$.
        Proving that there exists an oriented path $P''$ satisfying $\sigma^-(P'') \ge \sigma^-(P)$ can be done analogously.}
	\end{proof}

\end{document}
